\begin{table}[!h]
\centering
\caption{ERPT según el tamaño del shock cambiario}
\centering
\fontsize{8}{10}\selectfont
\begin{threeparttable}
\begin{tabular}[t]{cccccccc}
\toprule
$h$ & ERPT$^{-}$ (apr.) & ERPT$^{0}$ (norm.) & ERPT$^{+}$ (depr.) & $p$-valor & $F^{-}$ & $F^{0}$ & $F^{+}$\\
\midrule
0 & 0.005 [-0.006, 0.016] & -0.004 [-0.018, 0.011] & 0.000 [-0.010, 0.010] & 0.627 & --- & --- & ---\\
3 & 0.004 [-0.049, 0.057] & 0.002 [-0.036, 0.040] & 0.005 [-0.036, 0.047] & 0.993 & 22.6 & 28.2 & 54.1\\
6 & 0.015 [-0.068, 0.097] & 0.021 [-0.017, 0.058] & 0.057 [-0.054, 0.168] & 0.845 & 11.1 & 18.3 & 45.4\\
12 & 0.024 [-0.131, 0.179] & 0.045 [0.013, 0.077] & 0.234 [-0.339, 0.808] & 0.514 & 4.0 & 19.3 & 2.6\\
18 & 0.068 [-0.129, 0.264] & 0.067 [0.011, 0.123] & 0.256 [-0.369, 0.881] & 0.487 & 2.5 & 17.1 & 1.5\\
\addlinespace
24 & 0.085 [-0.066, 0.237] & 0.113 [0.031, 0.195] & 0.323 [-0.828, 1.474] & 0.194 & 3.9 & 17.3 & 0.3\\
\bottomrule
\end{tabular}
\begin{tablenotes}
\item \textit{Note: } 
\item Regímenes definidos por el umbral de $pm 1$ desvío estándar del shock. Intervalos al 95% con errores de Newey-West y $h+1$ rezagos. El $p$-valor corresponde al test de Wald de igualdad de los tres coeficientes de respuesta del IPC. $F^{-}$, $F^{0}$ y $F^{+}$ son los estadísticos de primera etapa de los regímenes de apreciación grande, normal y depreciación grande, respectivamente; el umbral convencional para descartar instrumento débil es 10.
\end{tablenotes}
\end{threeparttable}
\end{table}